\documentclass[12pt,a4paper]{article}
\usepackage[UTF8]{ctex}
\usepackage{geometry}
\usepackage{listings}
\usepackage{xcolor}
\usepackage{hyperref}

\geometry{margin=2.5cm}

\lstset{
    language=Python,
    basicstyle=\ttfamily\small,
    keywordstyle=\color{blue}\bfseries,
    commentstyle=\color{green!60!black},
    stringstyle=\color{red},
    numbers=left,
    numberstyle=\tiny\color{gray},
    stepnumber=1,
    numbersep=5pt,
    frame=single,
    breaklines=true,
    showstringspaces=false,
    tabsize=4
}

\title{Halo 系统代码文档：Util 模块}
\author{代码分析文档}
\date{\today}

\begin{document}

\maketitle

\tableofcontents
\newpage

\section{概述}

\texttt{util.py} 模块提供了一些通用的工具函数，主要用于：
\begin{itemize}
    \item 数据类型转换（字符串到 torch.dtype）
    \item GPU 设备管理（解析可见 GPU ID）
\end{itemize}

\section{函数详解}

\subsection{\_resolve\_dtype}

\subsubsection{功能}

将字符串类型的数据类型规范转换为 \texttt{torch.dtype} 对象。

\subsubsection{实现代码}

\begin{lstlisting}
def _resolve_dtype(dtype_spec):
    """
    Map common dtype strings (e.g., 'bfloat16', 'float16', 'float32') to torch dtypes.
    If dtype_spec is already a torch.dtype, return it directly.
    """
    if isinstance(dtype_spec, torch.dtype):
        return dtype_spec
    if isinstance(dtype_spec, str):
        key = dtype_spec.strip().lower()
        table = {
            "bf16": torch.bfloat16,
            "bfloat16": torch.bfloat16,
            "fp16": torch.float16,
            "float16": torch.float16,
            "f16": torch.float16,
            "float32": torch.float32,
            "fp32": torch.float32,
            "f32": torch.float32,
            "float": torch.float32,
            "half": torch.float16,
        }
        return table.get(key, torch.float32)
    # Fallback
    return torch.float32
\end{lstlisting}

\subsubsection{功能说明}

\begin{description}
    \item[输入] 可以是字符串或 \texttt{torch.dtype} 对象
    \item[处理]
    \begin{enumerate}
        \item 如果已经是 \texttt{torch.dtype}，直接返回
        \item 如果是字符串，转换为小写并查找映射表
        \item 如果找不到匹配，默认返回 \texttt{torch.float32}
    \end{enumerate}
    \item[支持的格式]
    \begin{itemize}
        \item \texttt{"bfloat16"} / \texttt{"bf16"}
        \item \texttt{"float16"} / \texttt{"fp16"} / \texttt{"f16"} / \texttt{"half"}
        \item \texttt{"float32"} / \texttt{"fp32"} / \texttt{"f32"} / \texttt{"float"}
    \end{itemize}
\end{description}

\subsubsection{使用场景}

\begin{itemize}
    \item 在加载模型时，将配置中的字符串类型转换为 PyTorch 类型
    \item 提供灵活的输入格式，支持多种常见的类型表示方式
\end{itemize}

\subsection{\_visible\_physical\_gpu\_ids}

\subsubsection{功能}

解析当前可见的物理 GPU ID 列表，考虑 \texttt{CUDA\_VISIBLE\_DEVICES} 环境变量的影响。

\subsubsection{实现代码}

\begin{lstlisting}
def _visible_physical_gpu_ids() -> list[int]:
    """
    Resolve the list of physical GPU IDs from CUDA_VISIBLE_DEVICES.
    If not set, fallback to [0, 1, ..., torch.cuda.device_count()-1].
    """
    env = os.environ.get("CUDA_VISIBLE_DEVICES", "").strip()
    if not env:
        return list(range(torch.cuda.device_count()))
    return [int(x) for x in env.split(",") if x.strip() != ""]
\end{lstlisting}

\subsubsection{功能说明}

\begin{description}
    \item[逻辑]
    \begin{enumerate}
        \item 检查环境变量 \texttt{CUDA\_VISIBLE\_DEVICES}
        \item 如果未设置，返回所有可用 GPU 的 ID（0 到 device\_count-1）
        \item 如果已设置，解析逗号分隔的 ID 列表
    \end{enumerate}
    \item[返回值] 物理 GPU ID 的整数列表
\end{description}

\subsubsection{使用场景}

\begin{itemize}
    \item 在创建 Worker 进程时，确定应该使用哪些 GPU
    \item 支持通过环境变量限制可用的 GPU
    \item 在多 GPU 系统中，可以只使用部分 GPU 进行测试
\end{itemize}

\subsubsection{示例}

\begin{lstlisting}
# 环境变量未设置，系统有 4 个 GPU
# 返回: [0, 1, 2, 3]

# 设置 CUDA_VISIBLE_DEVICES=0,2
# 返回: [0, 2]

# 设置 CUDA_VISIBLE_DEVICES=1
# 返回: [1]
\end{lstlisting}

\section{总结}

\texttt{util.py} 模块虽然简单，但提供了重要的基础设施：

\begin{itemize}
    \item \textbf{类型转换}：统一处理数据类型字符串，提高配置的灵活性
    \item \textbf{设备管理}：解析 GPU 可见性，支持灵活的 GPU 分配策略
\end{itemize}

这些工具函数被 Optimizer 和 Worker 广泛使用，是系统正常运行的基础。

\end{document}
